\chapter{Giới thiệu đề tài}
\label{Chapter1}

%Tóm tắt luận văn được trình bày nhiều nhất trong 24 trang in trên hai mặt giấy, cỡ chữ Times New Roman 11 của hệ soạn thảo Winword hoặc phần mềm soạn thảo Latex đối với các chuyên ngành thuộc ngành Toán.

%Mật độ chữ bình thường, không được nén hoặc kéo dãn khoảng cách giữa các chữ.
%Chế độ dãn dòng là Exactly 17pt.
%Lề trên, lề dưới, lề trái, lề phải đều là 1.5 cm.
%Các bảng biểu trình bày theo chiều ngang khổ giấy thì đầu bảng là lề trái của trang.
%Tóm tắt luận án phải phản ảnh trung thực kết cấu, bố cục và nội dung của luận án, phải ghi đầy đủ toàn văn kết luận của luận án.
%Mẫu trình bày trang bìa của tóm tắt luận văn (phụ lục 1).
\pagebreak
Phát hiện bất thường chuỗi thời gian (TSAD - Time Series Anomaly Detection) là một bài toán có nhiều ứng dụng trong các lĩnh vực quan trọng như sản xuất chất bán dẫn \cite{11456978}, giám sát hạ tầng điện toán đám mây \cite{10.1145/3637528.3671936} \cite{9966938} \cite{11121677}, điều trị nội trú và ngoại trú trong lĩnh vực y tế \cite{NEURIPS2023_7a53bf4e}, và nhiều ứng dụng khác. Các phương pháp TSAD thường dựa trên hai tác vụ học \textit{(learning task)} chính là \textit{dự báo (forecasting)} và \textit{tái tạo chuỗi (reconstruction)}. Trong đó, các phương pháp TSAD \textit{dựa trên tác vụ tái tạo chuỗi (reconstruction-based)} thường được sử dụng để phát hiện bất thường trong các lĩnh vực cần độ chính xác cao, do mô hình trong các phương pháp này dự đoán dựa trên dữ liệu mới nhất, không giống như các phương pháp \textit{dựa trên tác vụ dự báo (forecasting-based)} thường dựa trên dữ liệu lịch sử để dự đoán dữ liệu mới nhất. Trọng tâm của đề tài này là các phương pháp TSAD dựa trên tác vụ tái tạo chuỗi. Sau đây, chúng tôi sẽ thảo luận về khía cạnh dữ liệu của bài toán TSAD.

Gán nhãn bất thường cho dữ liệu chuỗi thời gian là một quá trình \textit{rất tốn nhân công (labor-intensive)}. Tuy nhiên, kể cả dữ liệu chuỗi thời gian đã được gán nhãn bình thường cũng luôn có khả năng bị \textit{nhiễm bất thường}. Hai lí do trên là động lực chính để một số phương pháp hiện nay sử dụng các kỹ thuật \textit{gia tăng dữ liệu (data-augmentation)} để thu được các mẫu, với một mẫu thường là một \textit{cửa sổ trượt}), có chứa \textit{bất thường nhân tạo (synthetic anomaly)} \cite{10.1145/3711896.3737110} \cite{jeong2026pgrfnet} \cite{darban_carla_2025}. Bằng cách này, mô hình sẽ có nhiều dữ liệu hơn để học cách trích xuất các mẫu hình bất thường. Đồng thời, các bất thường nhân tạo có thể vô tình gần giống với các bất thường thật chưa được gán nhãn trong dữ liệu bình thường. Điều này giúp mô hình tách biệt được rõ hơn các mẫu hình bất thường và các mẫu hình thật sự bình thường, bỏ qua \textit{giả định rằng chuỗi huấn luyện hoàn toàn bình thường (normality assumption)}.

Hiện nay, các phương pháp gia tăng số lượng mẫu dữ liệu bất thường dựa trên một trong hai giả định phổ biến là: \textit{giả định nhị phân về bất thường (binary anomaly assumption)}, nghĩa là có một lớp bình thường và một lớp bất thường, và \textit{giả định đa lớp về bất thường (multi-class anomaly assumption)}, nghĩa là có một lớp bình thường và nhiều lớp bất thường. Hai \textit{tác vụ học (learning task)} tương ứng với hai giả định này là \textit{phân loại nhị phân (binary classification)} và \textit{phân loại đa lớp (multi-class classification)}.

Gần đây, các phương pháp TSAD sử dụng mạng nơ-ron sâu (deep neural networks), bao gồm các phương pháp sử dụng kỹ thuật gia tăng dữ liệu, đã cho thấy kết quả khả quan. Tuy nhiên, việc định lượng độ bất định \textit{độ bất định của mô hình (epistemic uncertainty)} \cite{10.1145/3637528.3671936} trong tác vụ TSAD là một vấn đề chưa được quan tâm và nghiên cứu đủ nhiều. Những phương pháp không định lượng độ bất định của mô hình sẽ không nhận biết được \textit{mức độ tự tin (confidence)} và \textit{mức độ bất định (uncertainty)} của mô hình, đặc biệt khi sử dụng mô hình trên \textit{dòng dữ liệu chuỗi thời gian (time-series stream)} sinh ra từ phân phối khác với \textit{phân phối của tập huấn luyện (training distribution)}. Điều này làm giới hạn khả năng ứng dụng các phương pháp này trong các lĩnh vực rủi ro cao (high-stakes) như y tế kỹ thuật cao, tài chính - ngân hàng, và nhiều lĩnh vực khác.

Nhìn chung, các phương pháp hiện nay có hai hạn chế chung như sau: \newline
(1) Chưa lượng hoá được \textit{độ bất định (uncertainty)} trong kết quả phân lớp hay trong kết quả phát hiện bất thường. \newline
(2) Chưa thích ứng được mô hình một cách liên tục (\textit{continual adaptation}) và hiệu quả (\textit{efficiency}) trên các \textit{dòng dữ liệu thời gian có tính không dừng (non-stationary streaming time series)}.

Để giải quyết hai vấn đề trên, trong báo cáo này, chúng tôi đề xuất một phương pháp để định lượng độ bất định của mô hình thông qua sử dụng hàm Gumbel-Softmax \cite{jang2017categorical} trong cơ chế \textit{truy vấn nguyên mẫu (querying prototypes)}. Để đảm bảo thuật toán lan truyền ngược có thể tính toán gradient một cách đầy đủ, chúng tôi sử dụng \textit{kỹ thuật tái tham số hoá (re-parametrization technique)} để \textit{dòng gradient (gradient flow)} có thể đi qua hàm Gumbel-Softmax này và đến được các mô-đun trước đó của mạng nơ-ron.
Đồng thời, chúng tôi cũng đề xuất một phương pháp để thích ứng mô hình trên dòng dữ liệu thời gian có tính không dừng, với mục tiêu là đưa \textit{vec-tơ biểu diễn sâu tại thời điểm trực tuyến (online representation vector)} về gần với biểu diễn sâu mà mô hình học được sau khi trải qua quá trình \textit{tiền huấn luyện (pre-training)}, hay còn gọi là \textit{vec-tơ biểu diễn nguồn (source representation vector)}.
Ngoài ra, chúng tôi cũng sẽ đề xuất một phương pháp cải tiến biểu diễn mà mô hình học được, sao cho cho mỗi tác vụ trong \textit{khung huấn luyện đa tác vụ (multi-task learning framework)} mà chúng tôi sử dụng sẽ có biểu diễn chuyên biệt tương ứng, dựa trên việc sử dụng đồng thời: một nhánh truy vấn liên tục các vec-tơ nguyên mẫu (\textit{continuous query operator}) và một nhánh truy vấn rời rạc vec-tơ nguyên mẫu (\textit{discrete query operator}).

Tóm tắt lại thì trong báo cáo này chúng tôi sẽ đề xuất những điều sau: \newline
(1) Một toán tử truy vấn liên tục các vec-tơ nguyên mẫu và một toán tử truy vấn rời rạc các vec-tơ nguyên mẫu, kèm theo hai phiên bản \textit{truy vấn ngẫu nhiên (stochastic query)} tương ứng. \newline
(2) Một khung thích ứng trực tuyến để kéo vec-tơ biểu diễn trực tuyến về gần với vec-tơ biểu diễn gốc tương ứng, hay nói cách khác là đảm bảo \textit{bộ mã hoá trực tuyến (online encoder)} không trôi quá xa \textit{bộ mã hoá nguồn (source encoder)}. \newline
(3) Một cơ chế bộ đệm (verification buffer) để lọc ra các cửa sổ dùng cho \textit{quá trình cập nhật trọng số trực tuyến (online stochastic gradient descent)}. Điều kiện lọc sẽ được trình bày cụ thể trong \textit{Chương 3. Phương pháp đề xuất}.